\section{Integral topology and geometric observations}
\label{sec:integral-topology}

Throughout, a complete intersection has its scheme structure, and its
complex orientation fixes Poincar\'e duality. Put $A=\PP^3\times\PP^3$.
We first compute the integral topology of every smooth complete
intersection of bidegrees $(1,1),(1,1),(2,2)$. This computation does
not require its first two equations to cut out a smooth fourfold.

\begin{proposition}\label{prop:middle-lattice}
For any such smooth complete intersection $X$,
\[
 H_1(X;\ZZ)=0,\qquad H_2(X;\ZZ)=\ZZ^2,\qquad H_3(X;\ZZ)=\ZZ^{134}.
\]
The intersection form on $H_3(X;\ZZ)$ is unimodular and alternating.
\end{proposition}

\begin{proof}
Consider the affine coefficient space
\begin{equation}\label{eq:full-coefficients}
 \mathcal C=H^0(A,\OO(1,1))^{\oplus2}
                       \oplus H^0(A,\OO(2,2))\cong\CC^{132}.
\end{equation}
Its locus $\mathcal U$ of smooth three-dimensional intersections
is Zariski open: the rank-deficient zero locus of the universal
equations is closed and projective over $\mathcal C$, so its image
is closed. A regular triple has nonempty zero set because its
intersection class is nonzero. Successive Bertini for the very ample
bundles gives a triple whose entire successive flag is smooth.

The open set $\mathcal U$ is path connected. Two of its points
can be joined through a third point along complex lines not contained
in its algebraic complement; each line meets that complement in
finitely many points, which a path can avoid. The universal family
over $\mathcal U$ is proper and smooth. Transport along the resulting
path identifies the oriented integral topology of $X$ with that of
a smooth flag; see \cite[\S\S1,3]{Ehresmann} for the underlying
proper-submersion transport.

On the smooth flag the ambient dimensions are successively $6,5,4,3$.
Integral Lefschetz for a smooth very ample divisor therefore preserves
$H_1$ and $H_2$ at all three steps
\cite[Corollary~7.3 and Theorem~7.4]{Milnor}. Those groups of $A$
are respectively zero and $\ZZ^2$, proving the first two assertions.
It also proves connectedness. The deformation is used only to
compute topology; it does not replace the specific observation family.

All groups are finitely generated. Universal coefficients give
\[
 0\longrightarrow\operatorname{Ext}(H_2(X;\ZZ),\ZZ)
 \longrightarrow H^3(X;\ZZ)
 \longrightarrow\operatorname{Hom}(H_3(X;\ZZ),\ZZ)
 \longrightarrow0.
\]
The left term vanishes, so $H^3$ is free. Integral Poincar\'e duality
then makes $H_3$ free and its pairing unimodular
\cite[Theorems~3.2, 3.30 and Proposition~3.38]{Hatcher}.
Graded commutativity in degree three makes the pairing alternating.

For the rank, let $x,y$ be the ambient hyperplane classes and put
$s=x+y$. The tangent and normal bundle sequences give
\begin{equation}\label{eq:cicy-euler-data}
 [X]=2s^3,\qquad
 c(TX)=\frac{(1+x)^4(1+y)^4}{(1+s)^2(1+2s)},\qquad
 c_3(TX)=-2(x^3+y^3)-10(x^2y+xy^2).
\end{equation}
The coefficient of $x^3y^3$ in $2s^3c_3(TX)$ is $-128$.
Thus $\chi(X)=-128$. Connectedness and duality give
$b_0=b_6=1$, $b_1=b_5=0$, and $b_2=b_4=2$, whence $b_3=134$.
\end{proof}

\subsection{Markings and the observation map}

\begin{definition}\label{def:geometric-marking}
A geometric marking of $X$ consists of an integral basis
$\gamma_1,\ldots,\gamma_{134}$ of $H_3(X;\ZZ)$, finite
piecewise smooth integral singular cycles $c_a$ representing it,
and the intersection matrix $J_{ab}=\langle\gamma_a,\gamma_b\rangle$.
It is symplectic if $J$ is the chosen standard symplectic matrix.
\end{definition}

Such markings exist: a unimodular alternating lattice admits an
integral symplectic basis, and smooth singular homology represents
its classes by finite smooth chains. Section~\ref{sec:integral-marking}
constructs the marking and derivative bounds from algebraic fibre
equations. It does not require the coordinates of any vanishing cycle.

For a smooth proper family along a loop $\ell$, choose a smooth
horizontal splitting and let $F_\ell:X\longrightarrow X$ be the
endpoint transport. For a closed three-form $\Psi$ on $X$, define
\begin{equation}\label{eq:geometric-period-observations}
 \Pi_a=\int_{c_a}\Psi,\qquad
 \Pi_{a\ell}^{\rm end}=\int_{(F_\ell)_\#c_a}\Psi,\qquad
 \Delta_{a\ell}=\Pi_{a\ell}^{\rm end}-\Pi_a.
\end{equation}
These operations transport and integrate the actual chain. They
do not first decompose the endpoint in the original homology basis.
The latter coordinates, intersections with local vanishing spheres,
and monodromy matrices are not forward inputs.

\begin{lemma}\label{lem:transport-independence}
The quantities in \eqref{eq:geometric-period-observations} depend
only on the homology marking, the homotopy class of the loop with
fixed basepoint, and the cohomology class of $\Psi$. They are
independent of the horizontal splitting.
\end{lemma}

\begin{proof}
The pullback family over the compact path is proper. A horizontal
lift therefore exists for the whole interval and is invertible by
reversed transport. The right inverses defining two splittings are
joined in their affine space. Lifting this family gives an isotopy
between the endpoint diffeomorphisms. The same argument over a
homotopy gives homology continuation depending only on the path
class. Integration of a closed form depends only on homology.
\end{proof}

In our family $\Psi$ is a single-valued relative holomorphic top
form, hence closed on each smooth fibre. Its value at a loop endpoint
is the original $\Psi$, with no scalar monodromy to remove.
The geometric definition above is meaningful independently of
numerical access. The transport estimates below provide that access.
